\chapter{Fra\"{\i}ss\'{e} Theory}

We will review here some basic ideas of model theory concerning the Fra\"{\i}ss\'{e}construction and ultrahomogeneous countable structures.  

Our main reference
here is
\begin{itemize}
    \item Hodges \cite{H}, Ch. 7.
    \item See also \cite{Ch} and \cite{Ca}.
\end{itemize}



\begin{definition}
    A (countable) {\it signature} is a countable collection $L=\{R_i\}_{i\in I}  \cup\{f_j\}_{j\in J}$ of (distinct) {\it relation} and {\it function symbols}  each of which has an associated number, called its {\it arity}.

    \begin{itemize}
        \item The arity $n(i)$ of each relation symbol $R_i$ is a positive integer and the arity $m(j)$ of each function symbol $f_j$ is a non-negative integer.
        \item A structure
              for $L$ is an object of the form
              \[
                  \bfa=\langle A,\{R^\bfa_i\}_{i\in I},\{f^\bfa_j\}_{j\in J}\rangle ,
              \]
              where
              \begin{itemize}
                  \item $A$ is a {\it non-empty} set, called the {\it universe} of $\bfa$,
                  \item $R^\bfa_i\subseteq A^{n(i)}$, i.e., $R^\bfa_i$ is a $n(i)$-ary relation on $A$,
                  \item and $f_j:A^{m(j)}\rightarrow A$, i.e., $f_j^{\bfa}$ is an $m_j$-ary function on $A$. (When $m(j)=0,\ f_j^{\bfa}$ is a distinguished element of $A$.)
              \end{itemize}


    \end{itemize}


\end{definition}


\begin{definition}
    Given two structures $\bfa ,\bfb$ of the same signature $L$, a {\it
            homomorphism} of $\bfa$ to $\bfb$ is a map $\pi :A\rightarrow B$
    such that
    \[
        R^\bfa_i(a_1,\dots ,a_{n(i)})\iff R^\bfb_i(\pi (a_1),
        \dots ,\pi (a_{n(i)}))
    \]
    and
    \[
        \pi (f^\bfa_j(a_1,\dots ,a_{m(j)}))=f^\bfb_j(\pi (a_1),\dots ,\pi (a_{m(j)}
            )).
    \]
    We write also in this case $\pi :\bfa\rightarrow\bfb$. ({\it Caution.} Sometimes in the definition of homomorphism, one only requires the left-to-right implication concerning $R_i^{\bfa}, R_i^{\bfb}$. We will use here only the stronger version above.)

    \begin{itemize}
        \item If $\pi$ is also 1-1, it is called a {\it monomorphism} or {\it embedding}.
        \item if $\pi$ is 1-1 and onto it is called an {\it isomorphism}.
        \item If there  is an isomorphism from $\bfa$ to $\bfb$, we say that $\bfa ,\bfb$ are {\it isomorphic}, in symbols $\bfa\cong\bfb$.
        \item Finally, An {\it automorphism} of $\bfa$ is an isomorphism of $\bfa$ to itself.
        \item We denote by Aut$(\bfa )$  the group of automorphisms of $\bfa$.

    \end{itemize}

\end{definition}
We will be primarily interested in countable structures $\bfa$ (i.e., the
universe $A$ is countable).

\begin{theorem}
    For countable $\bfa$,
    \begin{itemize}
        \item the group Aut$(\bfa )$, with the pointwise convergence topology, is Polish, in fact it is a closed subgroup of $S_A$, the Polish group of permutations of $A$ with the pointwise convergence topology.

              \begin{mynote}
                  \begin{definition}
                      A topological space $X$ is \emph{Polish} if it is separable and completely metrizable; i.e., there is a metric $d$ such that $d$ induces the topology on $X$ and $(X,d)$ is complete.
                  \end{definition}

                  \begin{theorem}[Closed subspaces of Polish spaces are Polish]
                      Let $(X,d)$ be a Polish metric space and let $Y\subseteq X$ be closed (with the subspace topology). Then $Y$ is Polish.
                  \end{theorem}

                  \begin{proof}
                      We verify separability and complete metrizability for $Y$.

                      \smallskip
                      \noindent\textbf{(1) Completeness.}
                      Equip $Y$ with the restricted metric $d_Y:=d\!\upharpoonright_{Y\times Y}$. If $(y_n)$ is Cauchy in $(Y,d_Y)$, then it is Cauchy in $(X,d)$ and hence converges in $X$ to some $y\in X$. Since $Y$ is closed in $X$, necessarily $y\in Y$. Thus $(Y,d_Y)$ is complete.

                      \smallskip
                      \noindent\textbf{(2) Separability.}
                      Fix a countable dense set $D=\{x_k:k\in\Bbb N\}$ in $X$. For each $n\in\Bbb N$ consider the countable cover of $X$ by open balls $\{B(x_k,2^{-n}):k\in\Bbb N\}$. For every pair $(n,k)$ with $B(x_k,2^{-n})\cap Y\neq\emptyset$, choose and fix a witness point
                      \[
                          y_{n,k}\in Y\cap B(x_k,2^{-n}).
                      \]
                      Set
                      \[
                          E:=\bigcup_{n\in\Bbb N}\,\{\,y_{n,k}:\ B(x_k,2^{-n})\cap Y\neq\emptyset\,\}\ \subseteq Y.
                      \]
                      Then $E$ is a countable subset of $Y$. We claim $E$ is dense in $Y$. Let $U\subseteq Y$ be nonempty and open (in the subspace topology). Then $U=Y\cap V$ for some open $V\subseteq X$, and there exist $y\in U$ and $r>0$ with $B(y,r)\subseteq V$. Choose $n$ with $2^{-n}<r$ and choose $k$ with $y\in B(x_k,2^{-n})$ (possible because $D$ is dense). By construction $B(x_k,2^{-n})\cap Y\neq\emptyset$ and hence $y_{n,k}$ is defined; moreover
                      \[
                          y_{n,k}\in Y\cap B(x_k,2^{-n})\subseteq Y\cap B(y,r)\subseteq U.
                      \]
                      Thus every nonempty open $U\subseteq Y$ meets $E$, proving that $E$ is dense in $Y$.

                      \smallskip
                      Combining (1) and (2), $Y$ is separable and completely metrizable under $d_Y$, hence $Y$ is Polish.
                  \end{proof}

                  \begin{corollary}[Closed subgroups of Polish groups are Polish]
                      Let $G$ be a Polish topological group and let $H\leq G$ be a closed subgroup (with the subspace topology). Then $H$ is Polish. In particular, $H$ is a Polish topological group because the multiplication and inversion on $H$ are the restrictions of the continuous operations on $G$.
                  \end{corollary}

                  \begin{proof}
                      As a topological space, $H$ is a closed subspace of the Polish space $G$, hence Polish by the theorem. Continuity of the group operations on $H$ follows from the fact that they are restrictions of continuous maps on $G$.
                  \end{proof}

                  \begin{lemma}[Closed $\Rightarrow$ sequentially closed in metric spaces]
                      Let $(X,d)$ be a metric space and $Y\subseteq X$ closed. If $(y_n)$ is a sequence
                      in $Y$ with $y_n\to y$ in $X$, then $y\in Y$.
                  \end{lemma}

                  \begin{proof}
                      Suppose, toward a contradiction, that $y\notin Y$. Since $Y$ is closed,
                      $X\setminus Y$ is open, so there is $\varepsilon>0$ with
                      $B(y,\varepsilon)\subseteq X\setminus Y$. But $y_n\to y$ means there is
                      $N$ such that for all $n\ge N$, $d(y_n,y)<\varepsilon$, hence
                      $y_n\in B(y,\varepsilon)\subseteq X\setminus Y$, contradicting $y_n\in Y$.
                      Therefore $y\in Y$.
                  \end{proof}
              \end{mynote}

              %   \begin{mynote}
              %       I have its proof in my notes for "Automorphism groups of countable structures, David Evans". A minor difference would be that they have different topologies but:

              %       %---------------------------------------------
              %       % Definition: Pointwise-convergence topology
              %       %---------------------------------------------
              %       \begin{definition*}[Pointwise-convergence topology]
              %           Let $X$ be a set and $(Y,\tau_Y)$ a topological space. On the function set $Y^X$
              %           define the \emph{pointwise-convergence topology} as the \emph{initial topology}
              %           with respect to the evaluation maps
              %           \[
              %               \operatorname{ev}_x : Y^X \to Y,\qquad \operatorname{ev}_x(f):=f(x)\quad(x\in X).
              %           \]
              %           Equivalently, it is the coarsest topology on $Y^X$ making all $\operatorname{ev}_x$
              %           continuous. A subbasic open set has the form
              %           \[
              %               \operatorname{ev}_x^{-1}(U)=\{\,f\in Y^X : f(x)\in U\,\}\quad(x\in X,\ U\in\tau_Y).
              %           \]
              %       \end{definition*}

              %       \begin{remark*}
              %           If $Y$ is discrete (e.g.\ $Y=A$ a set with the discrete topology), then on $Y^X$
              %           the pointwise-convergence topology equals the product topology $\prod_{x\in X} Y$.
              %           In particular, for a subgroup $G\leq A^A$ and $g\in G$, a neighbourhood basis at $g$
              %           is given by the sets
              %           \[
              %               U(g;F):=\{\,h\in G : h(a)=g(a)\ \text{for all } a\in F\,\},\qquad F\subseteq X\ \text{finite}.
              %           \]
              %       \end{remark*}

              %       \begin{lemma*}\label{lem:topologies-coincide}
              %           Let $A=\{a_0,a_1,\dots\}$ be countable and discrete, put $S_A=\mathrm{Sym}(A)$ with the product topology,
              %           and let $G=\mathrm{Aut}(\mathbf A)\subseteq S_A$.
              %           The \emph{pointwise convergence} topology on $G$ agrees with the subspace topology inherited from $S_A$.
              %       \end{lemma*}

              %       \begin{proof}
              %           For $g\in G$ and a finite set $F\subseteq A$ define
              %           \[
              %               U(g;F)\;:=\;\bigl\{\,h\in G \mid h(a)=g(a)\text{ for all }a\in F\,\bigr\}.
              %           \]

              %           \smallskip
              %           \noindent\textbf{(1) Subspace topology.}
              %           In the product topology on $A^A$ (hence on $S_A$ and its subspace $G$), the sets
              %           \[
              %               V(g;F)\;:=\;\bigl\{\,h\in G \mid h(a)=g(a)\text{ for all }a\in F\,\bigr\}
              %           \]
              %           with finite $F$ form a neighbourhood base at $g$.
              %           But $V(g;F)=U(g;F)$ by definition, so the subspace topology on $G$ is generated by the family $\{U(g;F)\}$.

              %           \smallskip
              %           \noindent\textbf{(2) Pointwise‐convergence topology.}
              %           By definition, the pointwise topology on $G$ is exactly the topology generated by the same family $\{U(g;F)\}$.

              %           \medskip
              %           Since both topologies have the identical neighbourhood base at every point, they coincide.
              %       \end{proof}
              %   \end{mynote}
        \item Conversely, given a closed subgroup $G \subseteq S_A$, there is a signature $L$ and a structure $\bfa_G$ with universe $A$, so that Aut$(\bfa_G )=G$.

    \end{itemize}




\end{theorem}

\begin{proof}[Sketch Of Proof ]
    To see this, let for each $n\geq 1,\o_1^n,\o_2^n,\dots$ be the orbits of $G$ on $A^n$, the action of $G$ on $A^n$ being defined by $g\cdot (a_1,\dots ,a_n)=(g(a_1),\dots , g(a_n))$.  Let $L=\{R_{n,i}\}_{n\geq 1}$, where each $R_{n,i}$ is an $n$-ary relation symbol.  Define $\bfa_G$ by letting $R^{\bfa_G}_{n,i}=\o_i^n\subseteq A^n$.  Then it is easy to check that Aut$(\bfa_G)=G$.
\end{proof}

We call $\bfa_G$
the {\it induced structure} associated to $G$ (see Hodges \cite{H}, 4.1.4,
where this is called the canonical structure for $G$ - we use however the
term canonical for other purposes in this paper).


\bigskip
A substructure $\bfb$ of $\bfa$ has as universe a (non-empty) subset
$B\subseteq A$ closed under each $f^\bfa_j$, and $R^\bfb_i=R^\bfa_i
    \cap B^{n(i)},f^\bfb_j=f^\bfa_j|B^{m(j)}$. We write $\bfb\subseteq\bfa$ in this case. For each $X\subseteq A$, there
is a smallest substructure containing $X$, called the {\it substructure
        generated} by $X$.  A substructure is {\it finitely generated} it is
generated by a finite set.  A structure is {\it locally finite} if all its
finitely generated substructures are finite.  For example, if $L$ is
    {\it relational}, i.e., $J=\emptyset$, the substructure generated by $X$
has universe $X$ and so every finitely generated substructure is finite.
This is also true if $J$ is finite and each $f_j$ has arity 0.

A structure $\bfa$ is called {\it ultrahomogeneous} if every isomorphism
between finitely generated substructures $\bfb ,\bfc$ of $\bfa$ can be
extended to an automorphism of $\bfa$.  For example, $\langle\bbQ ,<
    \rangle$, the rationals with the usual order, is an ultrahomogeneous
structure.  Fraïssé's theory provides a general analysis of
ultrahomogeneous countable structures.

Let $\bfa$ be a structure for $L$.  The {\it age} of $A$, Age$(\bfa )$
is the collection of all finitely generated structures in $L$ that can
be embedded in $\bfa$, i.e., the closure under isomorphism of the
collection of finitely generated substructures of $\bfa$.  Clearly the class $\k =$ Age$(\a )$ is non-empty, and satisfies the following
two properties:

(i) {\it Hereditary property (HP)}:  If $\bfb\in\k$ and $\bfc$ is a finitely
generated structure that can be embedded in $\bfb$, then $\bfc\in\k$.

(ii) {\it Joint embedding property (JEP)}:  If $\bfb ,\bfc\in\k$, there
is $\bfd\in\k$ such that $\bfb ,\bfc$ can be embedded in $\bfd$.

When $\bfa$ is moreover ultrahomogeneous, it is easy to see that $\k =$
Age$(\bfa )$ satisfies also the following crucial property:

(iii) {\it Amalgamation property (AP)}:  If $\bfb ,\bfc ,\bfd\in\k$
and $f:\bfb\rightarrow\bfc ,g:\bfb\rightarrow\bfd$ are embeddings, then
there is $\bfe\in\k$ and embeddings $r:\bfc\rightarrow\bfe ,s:\bfd
    \rightarrow\bfe$ such that $r\circ f=s\circ g$.

We summarize:

\begin{proposition}
    \label{2.1}
    Let $\bfa$ be an ultrahomogeneous
    structure.  Then $\k =\;${\rm Age}$(\bfa )$ is non-empty,
    and satisfies HP, JEP and AP.
    \qed
\end{proposition}

If $\bfa$ is countable, then clearly {\rm Age}$(\bfa )$ contains only
countably many isomorphism types.  Abusing language, we say that a class
$\k$ of structures is {\it countable} if it contains only countably many
isomorphic types.

We now have the following main result of Fra\"{\i}ss\'{e} \cite{Fr}.

\begin{theorem}[Fra\"{\i}ss\'{e}]
    \label{2.2}
    Let $L$ be a signature and
    $\k$ a class of finitely generated structures for $L$, which is non-empty,
    countable, and satisfies HP, JEP and AP.
    Then there is a unique, up to isomorphism, countable structure $\bfa$
    such that $\bfa$ is ultrahomogeneous and $\k =$ {\rm Age}$(\bfa )$.
    \qed
\end{theorem}


We call this structure the {\it Fra\"{\i}ss\'{e} limit} of $\k$,
\[
    \bfa ={\rm Flim} (\k ).
\]
Thus a countable ultrahomogeneous structure is the Fra\"{\i}ss\'{e} limit of
its age.  For example, Age$(\langle\bbQ ,<\rangle )=$ the class of
finite linear orderings $=\l O$ and so $\langle\bbQ ,<\rangle =
    {\rm Flim} (\l O)$.

We note here the following alternative characterization of ultrahomogeneity.

\begin{proposition}
    \label{2.3}
    Let $\bfa$ be a countable structure.
    Then $\bfa$ is ultrahomogeneous iff it satisfies the following {\it
            extension property}:
    \par
    If $\bfb ,\bfc$ are finitely generated and can be embedded in $\bfa$, $f:\bfb\rightarrow\bfa ,g:\bfb
        \rightarrow\bfc$ are embeddings, then there is
    an embedding $h:\bfc\rightarrow\bfa$
    such that $h\circ g=f$.
    \qed
\end{proposition}

\begin{mynote}
    \[
        \begin{tikzcd}[column sep=large, row sep=large, ampersand replacement=\&]
            \bfb \arrow[r, "g"] \arrow[d, "f"'] \& \bfc \arrow[ld, dashed, "h"] \\
            \bfa
        \end{tikzcd}
    \]

    \begin{definition*}
        We begin by saying that a structure \( D \) is \emph{weakly homogeneous} if it has the property

        \begin{equation}
            \tag{1.7}
            \begin{aligned}
                \text{if } A, B & \text{ are finitely generated substructures of } D,\ A \subseteq B\ \text{ and } \\
                                & f\colon A \to D\ \text{is an embedding, then there is an embedding }             \\
                                & g\colon B \to D\ \text{which extends } f.
            \end{aligned}
        \end{equation}

        \[
            \begin{tikzcd}[ampersand replacement = \&]
                A \arrow[r, "f"] \arrow[d, hook] \& D \\
                B \arrow[ur, dashed, "g"']
            \end{tikzcd}
        \]

    \end{definition*}
\end{mynote}


It follows that if $\bfa$ is countable, $\bfb$ is countable ultrahomogeneous
with Age$(\bfa )\subseteq {\rm Age} (\bfb )$ and $\bfc\subseteq\bfa$ is
finitely generated, then every embedding $f:\bfc\rightarrow\bfb$ can be
extended to an embedding $g:{\bfa}\rightarrow {\bfb}$.

In particular, $\bfb$ is {\it universal} for the class of all countable
structures $\bfa$ whose age is contained in that of $\bfb$, i.e., every
such $\bfa$ can be embedded in $\bfb$.  For example, any countable linear
ordering can be embedded in $\langle\bbQ ,<\rangle$.

We will be primarily interested in the case when the classes $\k$ in \ref{2.2}
actually consist of {\it finite} structures and the Fra\"{\i}ss\'{e}
limit of $\k$ is countably infinite.  It will be convenient then to
introduce, for later use, the following terminology, where the {\it
        cardinality} of a structure $\bfa =\langle A,\dots\rangle$ is the
cardinality of its universe $A$.

\begin{definition}
    \label{2.4}
    Given a signature $L$, a {\rm Fra\"{\i}ss\'{e}
    class} in $L$ is a class of {\it finite} structures in $L$, which
    contains structures of arbitrary
    large (finite) cardinality, is countable, and satisfies HP, JEP and AP.  A
    {\rm Fra\"{\i}ss\'{e} structure} in $L$ is a countably {\it infinite} structure
    which is locally finite and ultrahomogeneous.
\end{definition}

Thus the map $\k\mapsto{\rm Flim} (\k )$ is a bijection between Fra\"{\i}ss\'{e}
classes and Fra\"{\i}ss\'{e} structures (up to isomorphism) with inverse the
map $\bfa\mapsto {\rm Age} (\bfa )$.

We would like to point out here that for $G$ a closed subgroup of $S_A$,
the induced structure $\bfa_G$ is ultrahomogeneous

\begin{mynote}
    Since each basic relation of $\bfa_G$ is a $G$-orbit, any isomorphism between finitely generated substructures preserves orbit membership and is therefore induced by some $g\in G$, which extends it to an automorphism of $\bfa_G$, so $\bfa_G$ is ultrahomogeneous.
\end{mynote}

and, since the
associated signature is relational, it is locally finite, so it is a Fra\"{\i}ss\'{e}
structure, provided $A$ is infinite.

Finally, for further reference, we recall the following definition. A
class $\k$ of structures for $L$ satisfies the {\it strong amalgamation
        property} (SAP) if for any $\bfa ,\bfb ,\bfc \in\k$ and embeddings $f:
    \bfa\rightarrow\bfb ,g:\bfa\rightarrow\bfc$, there is $\bfd\in\k$ and
embeddings $r:\bfb\rightarrow\bfd ,s:\bfc\rightarrow\bfd$ with $r\circ
    f=s\circ g$, such that moreover $r(B)\cap s(C)=r(f(A))\; (=s(g(A))$.

Similarly, we say that $\k$ satisfies the {\it strong joint embedding
        property} (SJEP) if for any $\bfa ,\bfb\in\k$ there is $\bfc\in\k$
and embeddings $f:\bfa\rightarrow\bfc ,g:\bfb\rightarrow\bfc$ such that
$f(A)\cap g(B)=\emptyset$.

\begin{remark}
    In retrospect, one can say that Fra\"{\i}ss\'e's construction
    was anticipated by Urysohn \cite{U}, who considered the special case of
    the class of finite metric spaces with rational distances.
    He constructed a countable metric space $\bfu_0$ (see Section 6, {\bf (E)} below)
    whose completion $\bfu$, known as the
        {\it Urysohn space}, is the unique universal ultrahomogeneous (with
    respect to isometries) complete separable metric space.  Note that we
    can view metric spaces $(X,d)$ as structures in a countable signature
    $L=\{R_q\}_{q\in\bbQ},\ R_q$ binary, identifying $(X,d)$ with $ \bfx
        =(X,R^\bfx_q)$, where $(x,y)\in R^\bfx_q\iff d(x,y)<q$.
\end{remark}
